\section{虚拟细胞是什么}
% 2
\begin{frame}[t]{虚拟细胞：输入、干预与响应}
\tagline{长期愿景覆盖多尺度生物状态；本次用转录组扰动响应预测作为具体例子。}
\centering
\begin{tikzpicture}
 \node[box,text width=3.0cm,minimum height=1.85cm] (a) {\ico{\large\faVials}\\[.08cm]\textbf{输入}\\[.1cm]\small 对照细胞的表达状态\\细胞类型 / 实验条件\\基因抑制或药物干预};
 \node[box,fill=PaleTeal,text width=2.75cm,minimum height=1.85cm,right=.62cm of a] (b) {\ico{\large\faProjectDiagram}\\[.08cm]\textbf{细胞预测模型}\\[.1cm]\small 学习背景与干预\\到响应的映射};
 \node[box,text width=3cm,minimum height=1.85cm,right=.62cm of b] (c) {\ico{\large\faChartBar}\\[.08cm]\textbf{输出}\\[.1cm]\small 干预后的表达变化\\或细胞群体分布\\候选基因与通路};
 \draw[flow] (a)--(b);\draw[flow] (b)--(c);
 \draw[flow,draw=Muted] (c.south)--++(0,-.6)-|node[pos=.25,below,font=\scriptsize,text=Muted]{\faFlask\enspace 真实实验检验与数据反馈}(a.south);
\end{tikzpicture}
\vspace{.30cm}\remark{细胞模型给出数值预测；预测是否可信，取决于任务定义、对照基线和评估方式。}
\source{概念示意，参考 [1,2,4,8]。通常比较两组细胞，而非同一细胞干预前后的配对轨迹。}
\end{frame}
